
% !TEX program = lualatex
\documentclass[11pt,a4paper]{ltjsarticle}
\usepackage{icat-common}

%-------------------------------------------------
\begin{document}

\title{顕現論における知覚原理}
\author{千葉将臣}
\date{2026-04-19}
\maketitle

\begin{abstract}
  本稿は、世界の顕現（Aletheia）を支える最下層の公理群である「知覚原理」を、幾何学的論証形式によって記述するものである。我々の知覚に内在する非対称性が、いかにしてフラクタル構造、四句分別、そして無明（Avidyā）という演算体系を必然的に要請するかを明らかにする。

  顕現論は、古典論理学など二元論における真理保存性が自己同一性を前提とすることを明示し、その前提をフラクタル性へと置き換えることで、二元論を特殊ケースとして包含する体系である。
\end{abstract}

\section*{序言}
世界はその全体性において知覚されることはない。知覚とは、本質的に「境界」の生成であり、高次元の「空」を低次元の「色」へと射影する限定的な演算プロセスである。この限定性こそが、静寂なる潜（Letheia）に亀裂を入れ、世界を顕（Aletheia）として立ち上がらせる根源的動因である。

\section{定義}

\begin{description}
  \item[定義 1.0] \textbf{世界}：知覚不可能性をその本質とする全一の状態。世界そのものは知覚されることはない。
  \item[定義 1.1] \textbf{ある（是）}：知覚によって捕捉可能な、低次元に射影された情報の形式を指す。
  \item[定義 1.2] \textbf{ない（非）}：知覚によって直接捕捉不可能であり、「ある」の欠如としてのみ認識される高次元の情報の形式を指す。「ない」の指し示しは「無記」と呼ぶ。
  \item[定義 1.3] \textbf{次元}：情報を記述・捕捉するために必要な自由度の尺度。本体系において「ない」は「ある」よりも高い次元を持つ。
  \item[定義 1.4] \textbf{境界}：知覚が「ある」と「ない」を区分する際に生じる不連続点。この再帰的な区分がフラクタルを構成する。
  \item[定義 1.5] \textbf{時間性}：時間性は「ない」が「ある」よりも高い次元を持つことに起因し、次元差を埋めるための境界生成プロセス（$\metanegation$）として発生する。
  \item[定義 1.6] \textbf{消去}：「ある」の知覚と「ない」の認識の同時成立により、「ない」の持つ高次元の可能性が低次元の「ある」として捕捉可能な形に収束する操作を消去と呼ぶ。
  \item[定義 1.7] \textbf{世界の立ち上がり}：世界が「ある」として知覚可能な状態へ移行する際、「ない」の高次元性と極限の確定が同時に生じる。この同時性が消去の発生を可能にする機序である。
\end{description}    
\section{公理}

\subsection{知覚公理}
\begin{description}
  \item[公理 1.1] 知覚は「ある」のみを捕捉し、「ない」を直接捉えることはできず、
  指し示しである「無記」のみが可能である。
  \item[公理 1.2] 世界は二分（区分）されるものとしてのみ知覚される。
  \item[公理 1.3] 「ない」は「ある」の欠如としてのみ知覚の遡及的対象となる。
\end{description}

すなわち、四句分別は「是・非」の二分がフラクタル的に一段階展開した必然的構造であり、四句はフラクタルな境界を記述する過不足ない最小の完備表現である。

\subsection{消去公理}
\begin{description}
  \item[公理 1.4] \textbf{消去の完了}：消去の収束は確定した値として完了する。
  この完了した値としてある概念の極限・無限が生ずる。
\end{description}

\section{定理}

\begin{description}
  \item[定理 1.1] \textbf{フラクタル構造の必然性}\\
  世界の認識において、知覚の再帰的分割（公理 1.2）が無限に繰り返されることにより、部分と全体の相似が生じ、世界はフラクタル構造として知覚されざるを得ない。

\paragraph{四句分別のフラクタル構造}
公理 1.2 による二分（是・非）は、定理 1.1 のフラクタル構造に従い、
以下のように一段階展開する。

\[
\underbrace{\text{是},\quad \text{非}}_{\text{是の区分（ある）・世俗諦}}
\qquad
\underbrace{\text{亦是亦非},\quad \text{非是非非}}_{\text{非の区分（ない）・勝義諦}}
\]
  
  \item[定理 1.2] \textbf{四句分別の完備性}\\
  「ある」「ない」の次元差（定義 1.1, 1.2）を保持したままフラクタル境界を過不足なく記述する最小の表現形式は、是、非、亦是亦非、非是非非の四句分別である。
  
  \item[定理 1.3] \textbf{生成と収束の知覚限界}\\
  フラクタル構造におけるプロセスの「生成」は知覚可能（色）であるが、その「収束（極限）」は知覚上の欠如（空）として現れる。
\end{description}

\section{補足}
二値論理（二元論）が世界を網羅的に記述できると考える誤謬は、「ある」と「ない」を同次元に配置し、「ない」を単なる反転可能な「ある」として扱うことによって生じる。これは情報の「次元」を無視し、真理を静的な自己同一性に閉じ込める誤謬(非中道の誤謬)である。

\section{要請}
本体系の演算（証明）を駆動させるためには、知覚の最小単位 $\eta$ が、フラクタル則（自己相似性、スケール不変性、収束則、非対称性）および閉包則（自己参照の閉じ損ね）を満たす「無明」として存在することを前提とする。

\section{付録}
本知覚原理に基づく収束演算は、常に一意な $0$ 次元不動点 $\dfix$ を指標として持つ。これは釈迦および龍樹が示した「空」の計算論的な実態である。

\end{document}
